\documentclass[compress, 9pt]{beamer}
%\documentclass[handout]{beamer} aurait supprim� les ``overlays''

\usetheme{}% le th�me choisi pour ma pr�sentation
%\useoutertheme{smoothbars}
%\useinnertheme{rounded}
\usepackage[french]{babel}
\usepackage[latin1]{inputenc}
\usepackage{amsmath}
\usepackage{amssymb}
\usepackage{geometry}
\usepackage{graphicx}%pour les includegraphics
\usepackage{array}% pour des tableaux particuliers. N�cessaire pour les r�v�ler colonne par colonne
\usepackage{pstricks}
\usepackage{epsfig}
\usepackage{pst-plot}

\setbeamertemplate{navigation symbols}{}





%\addtobeamertemplate{footline}{\insertframenumber/\inserttotalframenumber}
%pour ajouter un num�ro de page en bas.% une autre solution : \useoutertheme{shadow}
%\usepackage{beamerthemesplit}




\begin{document}


\frame{


\begin{pspicture}(0,0)(11,8)



\psline[linewidth=0.05cm]{->}(0,0.5)(11,0.5)
\rput[B](5.5,0.1){\Large $x_1$}
\psline[linewidth=0.05cm]{->}(0.5,0)(0.5,7.5)
\rput[B](0.2,4){\Large $x_2$}

\psline[linewidth=0.05cm]{<-}(0,7)(11,7)
\rput[B](5.5,7.3){\Large $y_1$}
\psline[linewidth=0.05cm]{<-}(10.5,0)(10.5,7.5)
\rput[B](10.8,4){\Large $y_2$}


\rput(0,0){\psframebox[linewidth=1pt,linecolor=red]{\LARGE  \red \bf $A$}}
\rput(11,7.5){\psframebox[linewidth=1pt,linecolor=blue]{\LARGE  \blue \bf $B$}}



\psplot[algebraic,linecolor=red,linestyle=solid,linewidth=0.05cm]{1}{6}{0.5+ 3/(x-0.5)}
\psplot[algebraic,linecolor=red,linestyle=dotted,linewidth=0.05cm]{2}{7}{0.5+ 8/(x-0.5)}
\psline[linewidth=0.03cm,linecolor=red,linestyle=dashed]{->}(2.5,2.5)(4,4)
\rput[B](4.5,4.2){\Large \red  $\nearrow$ in $A$'s utility }


\psplot[algebraic,linecolor=blue,linestyle=solid,linewidth=0.05cm]{10}{5}{7- 3/(10.5-x)}
\psplot[algebraic,linecolor=blue,linestyle=dotted,linewidth=0.05cm]{9}{4}{7- 8/(10.5-x)}
\psline[linewidth=0.03cm,linecolor=blue,linestyle=dashed]{->}(8.5,5)(7,3.5)
\rput[B](6.7,3.1){\Large \blue  $\nearrow$ in $B$'s utility }




\end{pspicture}


}


























\end{document}
